\subsection{Neutron Stars: Exploratory Data Analysis}
\label{sec:eda-ns}
\input{41_neutron_stars_numbers}

After filtering (Section~\ref{sec:neutron-stars}), the neutron-star table holds $\nsEdaRows$ configurations on $\nsEdaCurves$ curves, one per $(\beta, \lambda)$ pair. The analysis leads to five decisions, each argued below with its evidence:
\begin{enumerate}
  \item the central density enters as $\log_{10}\rhoc$, standardized;
  \item the charge target is $\log_{10}(\Dch/M)$, not $\log_{10}\Dch$;
  \item both free parameters, $\beta$ and $\lambda$, stay inputs;
  \item the surrogate is valid only on the triangular $(\beta, \lambda)$ region and the sampled densities, where the data pile up near the maximum mass;
  \item the split groups whole curves: GroupKFold over $(\beta, \lambda)$ for interpolation, with the rim held out to measure extrapolation.
\end{enumerate}

\paragraph{1.\ The central density enters as $\log_{10}\rhoc$, standardized.}
The density is positive and multiplicative, and the mass follows $\log_{10}\rhoc$ as one smooth S-shaped curve per $(\beta, \lambda)$ pair (\cref{fig:ns-mass-density}); the input is standardized to zero mean and unit variance on the training folds only. All curves share one grid of $\nsEdaDensityGridValues$ densities over $\nsEdaRhocOrders$ orders of magnitude, spaced about evenly in $\rhoc$, so the marginal is a plateau that falls off as curves end at their maximum mass (\cref{fig:ns-univariate-continuous}). The skew drops from $\nsEdaRhocRawSkew$ to $\nsEdaRhocLogSkew$ under the logarithm (\cref{tab:ns-univariate}), but given that shape it is weak evidence and not the argument.
\todo{Methodology: compare raw against $\log_{10}\rhoc$ inputs by model error.}

\InputIfFileExists{41_neutron_stars_fig_univariate_continuous}{}{}
\InputIfFileExists{41_neutron_stars_fig_mass_density}{}{}
\InputIfFileExists{41_neutron_stars_tab_univariate}{}{}
\FloatBarrier

\paragraph{2.\ The charge target is $\log_{10}(\Dch/M)$.}
The charge spans $\nsEdaDOrders$ orders of magnitude, so its target is logarithmic; dividing by the mass makes it a better one (\cref{tab:ns-charge-correlation}, \cref{fig:ns-charge-target}). Two coefficients measure how a target depends on another quantity. Pearson's $r$ measures how closely the points follow a straight line, and $r^2$ is the share of the target's variance that a straight line in that quantity explains. Spearman's $\rho$ is Pearson's $r$ computed on ranks: it asks only whether a larger value of one goes with a larger value of the other, whatever the shape of the relation.
Against the mass, $\log_{10}\Dch$ has $\rho = \nsEdaCorrMLogD$: heavier stars tend to carry more charge, so part of the charge merely repeats the mass. For $\log_{10}(\Dch/M)$, $\rho$ falls to $\nsEdaCorrMLogDM$ (and $r$ from $\nsEdaPearsonMLogD$ to $\nsEdaPearsonMLogDM$): ranking the stars by mass says almost nothing about their normalized charge, so the charge output no longer duplicates the mass output and the network learns the two separately.
Against $\lambda$, $r$ rises from $\nsEdaCorrLambdaLogD$ to $\nsEdaCorrLambdaLogDM$: a straight line in $\lambda$ explains $\nsEdaLambdaLinearPercentLogDM\%$ of the variance of $\log_{10}(\Dch/M)$ against $\nsEdaLambdaLinearPercentLogD\%$ of $\log_{10}\Dch$, so the normalized charge is more directly a function of the theory's parameter.
Against $\beta$, both targets stay near $r = \nsEdaCorrBetaLogDM$: the coupling has no overall straight-line effect, yet at fixed $\lambda$ it explains much of the variance (decision~3), so its influence depends on $\lambda$, an interaction the network has to capture.
\todo{Review: physical reading of $\Dch/M$ as the dimensionless charge-to-mass ratio.}

\InputIfFileExists{41_neutron_stars_tab_charge_correlation}{}{}
\InputIfFileExists{41_neutron_stars_fig_charge_target}{}{}
\FloatBarrier

\paragraph{3.\ Both free parameters stay inputs.}
$\lambda$ carries the strongest linear signal in the charge, and at fixed $\lambda$ the coupling $\beta$ still explains $\nsEdaBetaSharePercent\%$ of the variance of $\log_{10}(\Dch/M)$ (median over $\lambda$).
\todo{Review: $\beta$ as the first-order coupling and $\lambda$ as a second-order self-interaction, against the field equations (Section~\ref{sec:action}).}

\paragraph{4.\ The domain is triangular, and the data pile up near the maximum mass.}
Only $\nsEdaGridFillPercent\%$ of the $\nsEdaGridCells$ $(\beta, \lambda)$ cells hold a converged curve: the corner of large $\beta$ and large $\lambda$ is empty, and the maximum mass falls towards it, from $\nsEdaMassMax\,\Msun$ at small $\lambda$ (\cref{fig:ns-mass-max}, right). Curves also thin out towards that edge, from $\nsEdaRowsPerCurveMax$ rows to $\nsEdaRowsPerCurveMin$. Along each curve the densities are spaced about evenly while $dM/d\rhoc \to 0$ at the turning point, so $\nsEdaPileUpPercent\%$ of all rows lie within $\nsEdaPileUpWindow\,\Msun$ of their curve's maximum (\cref{fig:ns-mass-max}, left): a sampling artefact that weights the region of Hypothesis~$H_2$ heavily.
\todo{Review: $\nsEdaMassMax\,\Msun$ against the heaviest observed neutron stars (about $2.08\,\Msun$, PSR J0740+6620; about $2.35\,\Msun$, PSR J0952$-$0607) and theoretical maxima of about $2.2$ to $2.5\,\Msun$; the single polytropic equation of state sets the dataset's maximum; masses assumed in solar units; add citations.}

\InputIfFileExists{41_neutron_stars_fig_mass_max}{}{}
\FloatBarrier

\paragraph{5.\ The split groups whole curves.}
The rows of a curve are each other's nearest neighbours: for $\nsEdaSameCurvePercent\%$ of the rows the nearest row in standardized $(\beta, \lambda, \log_{10}\rhoc)$ lies on the same curve, at a median distance of $\nsEdaWithinMedian$ against $\nsEdaAcrossMedian$ to another curve. A random split of rows therefore leaks; a one-nearest-neighbour probe on $\log_{10}(\Dch/M)$ measures how much (\cref{tab:ns-split-strategies}, $\nsEdaFolds$ folds). Holding out rows at random gives an error of $\nsEdaSplitRandomRowsDM$, holding out whole curves $\nsEdaSplitCurvesDM$, the honest test of prediction at unseen $(\beta, \lambda)$ pairs. Whole $\beta$ lines cost the same ($\nsEdaSplitBetaLinesDM$), whole $\lambda$ lines are harder ($\nsEdaSplitLambdaLinesDM$), and the rim predicted from the interior is hardest ($\nsEdaSplitRimDM$), so extrapolation is reported separately and not claimed. For the mass the order reverses ($\nsEdaSplitRandomRowsM$ against $\nsEdaSplitCurvesM$): the mass depends mostly on $\rhoc$, and a neighbouring curve at the same density predicts it better than the next density on the same curve.

\InputIfFileExists{41_neutron_stars_tab_split_strategies}{}{}
\FloatBarrier
